\input{papers/template/style.tex}

\begin{document}
  \title{From Papers to Causal Models: Extracting, Validating, and
    Executing Financial Theories}

  \author{ADIA Authors}
  % TODO(*): Fill in the affiliation.
  \address{\textless{}Department\textgreater{},
    \textless{}Organization\textgreater{},
    \textless{}City, State/Country\textgreater{}}
  \email{@adia.ae}

  \keywords{causal inference, factor investing, large language models,
    structural causal models, paper-to-code}

  \date{\today}

  \begin{abstract}
    Published finance research rarely states the causal graph behind its
    empirical claims. Factor-investing studies in particular justify model
    specifications through correlations, which can produce statistically
    valid but causally misspecified models (``factor mirages''). Manual
    replication and causal re-analysis of the literature do not scale.

    This paper proposes a system that converts a finance paper into an
    executable theory. The system extracts variables, causal claims, and the
    empirical specification into a typed intermediate representation (the
    Theory Spec), validates the implied causal directed acyclic graph (DAG)
    against the authors' own specification, and generates three artifacts: a
    simulator, a replication module, and a causal estimator. Generated
    outputs are verified against the tables reported in the source paper.

    The paper presents the design and an evaluation protocol. No empirical
    evaluation has been carried out yet.
  \end{abstract}

  \maketitle

  \setcounter{tocdepth}{2}
  \tableofcontents

  % Status: OUTLINE stage. Each section lists the content it will contain,
  % derived from `notes.md`. Bullets are replaced by prose after the outline
  % is approved.

  \section{Introduction}\label{sec:introduction}

  \begin{itemize}
  \tightlist
  \item
    Problem: finance papers state empirical associations, while the
    underlying causal theory is implicit and unverified.

    \begin{itemize}
    \tightlist
    \item
      Most factor-investing articles make associational claims: no causal
      graph, specifications justified by correlations, no falsification
      experiments \cite{lopezdeprado2023causal}.
    \item
      Consequence: the ``factor mirage'', i.e., a model that looks
      statistically valid but is causally misspecified
      \cite{lopezdeprado2025primer}.
    \item
      Replication evidence shows the cost of unverified claims
      \cite{harvey2016cross,hou2020replicating,jensen2023replication}.
    \end{itemize}
  \item
    Prior work falls in four categories, each with a gap:

    \begin{itemize}
    \tightlist
    \item
      Causal inference provides the formal language (SCMs, do-calculus,
      identification) but is not applied to published papers at scale.
    \item
      Causal factor investing proposes a protocol but applies it by hand.
    \item
      Causal extraction from text recovers graphs but stops before the
      estimator, and rarely targets finance.
    \item
      Paper-to-code agents generate code from papers but ignore causal
      validity.
    \end{itemize}
  \item
    This paper fills the gap: a pipeline from PDF to causally validated,
    executable model, built around a typed Theory Spec of which the DAG is
    one view.
  \item
    Simplifying assumptions (bullet list, bold lead-in each):

    \begin{itemize}
    \tightlist
    \item
      \textbf{Empirical papers first.} Equilibrium and theory papers are
      routed to a separate path (Section~\ref{sec:limitations}).
    \item
      \textbf{Cross-sectional anomalies as initial scope.} Small DAGs and
      public ground-truth implementations \cite{chen2022open}.
    \item
      \textbf{Public data only.} The simulator must run without
      proprietary data.
    \item
      \textbf{Firm-level nodes only.} Macroeconomic nodes, time-unrolling,
      and point-in-time macro data are treated in a companion paper
      \cite{adia2026macro}.
    \end{itemize}
  \item
    Contributions, each tied to a section:

    \begin{itemize}
    \tightlist
    \item
      The Theory Spec intermediate representation (Section~\ref{sec:theory_spec}).
    \item
      A six-stage pipeline with causal validation and feedback routing
      (Section~\ref{sec:architecture}).
    \item
      An automated ``factor mirage'' diagnostic: theory DAG vs.~empirical
      specification (Section~\ref{sec:architecture}).
    \item
      An evaluation protocol for the extraction and code halves (Section~\ref{sec:evaluation}).
    \end{itemize}
  \item
    Roadmap: Section~\ref{sec:related_work} reviews related work,
    Section~\ref{sec:problem} formulates the problems, and
    Section~\ref{sec:theory_spec} defines the Theory Spec.
    Section~\ref{sec:architecture} describes the pipeline, and
    Section~\ref{sec:evaluation} presents the evaluation protocol.
    Section~\ref{sec:limitations} discusses limitations and
    Section~\ref{sec:conclusion} concludes.
  \end{itemize}

  \section{Related Work}\label{sec:related_work}

  \begin{itemize}
  \tightlist
  \item
    \textbf{Causal inference foundations.} SCMs, do-calculus, backdoor and
    frontdoor identification \cite{pearl2009causality,pearl2018book};
    code-friendly treatment \cite{peters2017elements}; bridge between
    DAGs and potential outcomes \cite{imbens2020potential}; econometric
    designs (IV, difference in differences) as graphical identification
    \cite{hunermund2019causal}. Contrast: these supply the language, not
    the extraction or the execution.
  \item
    \textbf{Causality in finance.} Causal factor investing
    \cite{lopezdeprado2023causal,lopezdeprado2024case}, the primer and
    seven-step protocol
    \cite{lopezdeprado2025primer,lopezdeprado2025protocol}; replication
    literature
    \cite{harvey2016cross,hou2020replicating,jensen2023replication}.
    Contrast: the protocol is manual. This paper automates it and uses the
    protocol as the validator checklist.
  \item
    \textbf{Causal structure from text.} Pre-LLM survey
    \cite{yang2022survey}; FinCausal shared tasks; LLM causal reasoning
    and its failure modes
    \cite{kiciman2023causal,jin2023corr2cause,zecevic2023causal}; graph
    extraction by iterated pairwise queries and by retrieval with
    self-consistency \cite{zeroshot2023causal,lacr2024}; ReCast
    benchmark of expert-annotated economics graphs \cite{recast2025}.
    Contrast: no work carries the extracted graph through identification
    and into code.
  \item
    \textbf{Paper-to-code and LLM quant agents.} PaperCoder's plan,
    analyze, generate split \cite{seo2026paper2code}; rubric-based
    replication grading \cite{starace2025paperbench}; RD-Agent(Q)
    \cite{li2025rdagent}; FactorEngine \cite{factorengine};
    open-source signals \cite{chen2022open}. Contrast: these generate
    code without checking that the estimand is identified from the paper's
    own graph.
  \end{itemize}

  \section{Problem Formulation}\label{sec:problem}

  \subsection{Notation and Setup}\label{sec:notation_and_setup}

  \begin{itemize}
  \tightlist
  \item
    A source paper \(P\); a Theory Spec
    \(\mathcal{T} = (G, \mathcal{A}, \tau,
    \mathcal{S})\) with DAG \(G = (V, E)\), assumptions \(\mathcal{A}\),
    estimand \(\tau\), and empirical specification \(\mathcal{S}\).
  \item
    Node attributes: observed or latent, unit, time index, level (firm,
    industry).
  \item
    Edge attributes: sign, lag, status (\texttt{stated} or
    \texttt{inferred}), provenance.
  \end{itemize}

  \subsection{Simplifying Assumptions}\label{sec:simplifying_assumptions}

  \begin{itemize}
  \tightlist
  \item
    Bullet list, bold lead-in each: time-indexing removes feedback cycles;
    availability timing is encoded in the Theory Spec; the finance
    ontology is small and fixed; variable recipes can be expressed with
    standard data sources.
  \end{itemize}

  \subsection{Problem 1: Extraction}\label{sec:problem_1_extraction}

  \begin{itemize}
  \tightlist
  \item
    Map \(P\) to \(\mathcal{T}\) such that every node, edge, and
    assumption carries a provenance pointer. The pointer is
    format-specific: page and span (PDF), file, line, and label (LaTeX),
    heading path and line (Markdown).
  \item
    The input \(P\) may be a PDF, a LaTeX source bundle, or a Markdown
    file.
  \item
    Quality measured by structural distance and edge-level precision and
    recall against expert graphs.
  \end{itemize}

  \subsection{Problem 2: Causal
  Validation}\label{sec:problem_2_causal_validation}

  \begin{itemize}
  \tightlist
  \item
    Decide whether \(\tau\) is identifiable from \(G\).
  \item
    Decide whether \(\mathcal{S}\) is consistent with \(G\): no collider
    or mediator controls, no omitted confounders.
  \item
    List the conditional independencies implied by \(G\) for testing.
  \end{itemize}

  \subsection{Problem 3: Executable Model
  Generation}\label{sec:problem_3_executable_model_generation}

  \begin{itemize}
  \tightlist
  \item
    Map \(\mathcal{T}\) to a simulator, a replication module, and a causal
    estimator.
  \item
    Verification: reproduced outputs match reported tables within a
    tolerance \(\epsilon\).
  \end{itemize}

  \subsection{Problem 4: DAG Optimization
  (open)}\label{sec:problem_4_dag_optimization_open}

  \begin{itemize}
  \tightlist
  \item
    Select a revised DAG \(G^{*}\) that minimizes an objective \(J(G)\).
  \item
    The criterion \(J\) is deliberately unspecified. Candidates:
    out-of-sample fit, number of untestable assumptions, agreement with
    the literature-level graph.
  \end{itemize}

  % TODO(*): Notes defer the optimization criterion. Decide on J.

  \section{The Theory Spec}\label{sec:theory_spec}

  \begin{itemize}
  \tightlist
  \item
    Why not DAG to code: a bare DAG lacks data recipes, estimators, and
    timing.
  \item
    Schema (Pydantic or JSON):

    \begin{itemize}
    \tightlist
    \item
      \texttt{nodes}: id, canonical name, observed or latent, data recipe,
      unit, time index
    \item
      \texttt{edges}: source, target, sign, lag, functional-form hint,
      confidence, provenance, \texttt{stated} vs.~\texttt{inferred}
    \item
      \texttt{assumptions}: exogeneity, no unobserved confounding
    \item
      \texttt{estimand}: e.g., ATE of a characteristic on next-month
      return
    \item
      \texttt{empirical\_spec}: regressions, sorts, controls, fixed
      effects, sample
    \end{itemize}
  \item
    Canonicalization against a small finance ontology (value, size,
    momentum, accruals, leverage) so graphs from different papers merge.
  \item
    Rule: generated code never depends silently on an \texttt{inferred}
    edge.
  \item
    Table~\ref{tab:theory_spec_fields} lists the schema fields and their role in later stages.
  \end{itemize}

  \begin{table}[H]
    \centering
    \begin{tabularx}{\linewidth}{lXX}
      \toprule
      Field & Content & Role in later stages \\
      \midrule
      \texttt{nodes} & Id, canonical name, observed or latent, data recipe,
      unit, time index & Data construction in code generation (Stage 5) and
      verification (Stage 6) \\
      \texttt{edges} & Source, target, sign, lag, functional-form hint,
      confidence, provenance, \texttt{stated} or \texttt{inferred} &
      Identification and independence tests (Stage 4), simulator structure
      (Stage 5) \\
      \texttt{assumptions} & Exogeneity, no unobserved confounding &
      Identification (Stage 4) \\
      \texttt{estimand} & E.g., ATE of a characteristic on next-month return
      & Identification (Stage 4), causal estimator (Stage 5) \\
      \texttt{empirical\_spec} & Regressions, sorts, controls, fixed effects,
      sample & Factor mirage diagnostic (Stage 4), replication module
      (Stage 5), verification (Stage 6) \\
      \bottomrule
    \end{tabularx}
    \caption{Theory Spec schema fields.}
    \label{tab:theory_spec_fields}
  \end{table}

  \section{System Architecture}\label{sec:architecture}

  \begin{figure}[H]
    \centering
    % TODO(*): Replace the placeholder with the pipeline diagram.
    \fbox{\parbox{0.9\linewidth}{\centering\vspace{3em}
      Figure placeholder\vspace{3em}}}
    \caption{Pipeline from input document to verified executable model.}
    \label{fig:pipeline}
  \end{figure}

  \begin{itemize}
  \tightlist
  \item
    Figure~\ref{fig:pipeline} shows the pipeline: PDF, LaTeX, or Markdown
    input; Parse; Claim extraction; Theory Spec; Causal validation; Code
    generation; Execute and verify; feedback loop.
  \item
    \textbf{Stage 1, Parse.} One front end per input format, all producing
    the same normalized document representation (sections, equations,
    tables, bibliography, source anchors). Stages 2 to 6 read only this
    representation, so they are format-agnostic. Tables become
    verification targets.

    \begin{itemize}
    \tightlist
    \item
      \textbf{PDF.} Layout and OCR based extraction (Nougat, Marker, or
      GROBID). Fallback when no source is available. Equation and table
      errors are possible.
    \item
      \textbf{LaTeX.} Parse the source directly: flatten
      \texttt{\textbackslash{}input} and \texttt{\textbackslash{}include}
      files, expand custom macros, read \texttt{equation} and
      \texttt{tabular} environments, resolve
      \texttt{\textbackslash{}label}, \texttt{\textbackslash{}ref}, and
      \texttt{\textbackslash{}cite} keys, and read the \texttt{.bib} file.
      Equations and tables are exact. Preferred when the source is
      available (e.g., an arXiv source bundle).
    \item
      \textbf{Markdown.} Parse the AST (e.g., with pandoc): headings,
      \texttt{\$...\$} math, pipe tables, and citation keys. Quality
      depends on the converter when the Markdown was generated from
      another format.
    \item
      Format choice: LaTeX if available, then Markdown, then PDF.
    \item
      Open issue: tables generated by code and included as external files
      are absent from the source text.
    \end{itemize}
  \item
    \textbf{Stage 2, Claim extraction.} Multi-pass LLM extraction of
    variables, causal claims, empirical specification, and testable
    implications. Every item has a provenance pointer, as the main defense
    against hallucinated edges.
  \item
    \textbf{Stage 3, Theory Spec.} Assemble and canonicalize.
  \item
    \textbf{Stage 4, Causal validation.} DoWhy, pgmpy, or y0:

    \begin{itemize}
    \tightlist
    \item
      enforce acyclicity, handle feedback by time-indexing
    \item
      run identification of \(\tau\) from \(G\)
    \item
      compare the theory DAG with the empirical specification (the
      automated factor mirage diagnostic)
    \item
      list implied conditional independencies
    \end{itemize}
  \item
    \textbf{Stage 5, Code generation.} Plan, analyze, code. Three
    artifacts: simulator (SCM with known effects), replication module
    (data construction plus the paper's estimator), causal estimator
    (identified estimate). Boilerplate (CRSP and Compustat merges,
    portfolio sorts, Newey-West) is a tested library that the LLM composes
    rather than rewrites.
  \item
    \textbf{Stage 6, Execute and verify.} Sandbox run; compare to reported
    tables; route mismatches: wrong data recipe to Stage 2, missing
    confounder to Stage 3, bug to Stage 5.
  \end{itemize}

  \section{Evaluation Protocol}\label{sec:evaluation}

  \begin{itemize}
  \tightlist
  \item
    Status: protocol only. No results are reported.
  \item
    Initial scope: cross-sectional anomaly papers, hundreds of
    ground-truth implementations \cite{chen2022open}.
  \item
    \textbf{Extraction.} Normalized Structural Hamming Distance and
    edge-level precision and recall, ReCast style \cite{recast2025};
    per-node and per-edge LLM judging.
  \item
    \textbf{Code.} Correlation with published signals, matching
    t-statistics \cite{chen2022open}; rubric-based grading
    \cite{starace2025paperbench}.
  \item
    \textbf{Validation.} Share of papers with a flagged mirage.
  \item
    Baselines: direct paper-to-code without a DAG; pipeline without the
    validation stage.
  \item
    \textbf{Input format ablation.} Run the same papers from PDF, LaTeX,
    and Markdown inputs. Measure how parsing errors (equations, tables,
    citations) propagate to extraction and replication metrics.
  \item
    Table~\ref{tab:evaluation} lists metrics, data, and baselines.
  \end{itemize}

  \begin{table}[H]
    \centering
    \begin{tabularx}{\linewidth}{lXXX}
      \toprule
      Component & Metrics & Data & Baseline or variant \\
      \midrule
      Extraction & Normalized SHD, edge-level precision and recall, per-node
      and per-edge LLM judging & ReCast benchmark & Not applicable \\
      Code & Correlation with the published signal, matching t-statistics,
      rubric score & Chen--Zimmermann signals & Direct paper-to-code without a
      DAG \\
      Validation & Share of papers with a flagged mirage & Anomaly papers &
      Pipeline without the validation stage \\
      Input format & Propagation of parsing errors to extraction and
      replication metrics & The same papers as PDF, LaTeX, and Markdown & PDF
      front end \\
      \bottomrule
    \end{tabularx}
    \caption{Evaluation protocol: components, metrics, data, and baselines.}
    \label{tab:evaluation}
  \end{table}
  % TODO(*): Replace with a Technical Evaluation section once an
  % implementation exists.

  \section{Discussion and Limitations}\label{sec:limitations}

  \begin{itemize}
  \tightlist
  \item
    \textbf{Theory papers.} Equilibrium models (simultaneous equations,
    fixed points) do not map to a DAG. Route to a structural equations and
    solver path.
  \item
    \textbf{Implicit graphs.} Most papers never state a DAG.
    \texttt{stated} and \texttt{inferred} edges are kept separate. Quality
    of inferred edges depends on the LLM \cite{zecevic2023causal}.
  \item
    \textbf{Input availability.} Most finance papers circulate as PDF
    only. LaTeX sources exist mainly for preprints, and journal Markdown
    is rare, so the PDF front end remains the main path.
  \item
    \textbf{Look-ahead bias.} Reporting lags and fiscal-year alignment
    are the main failure points of generated code.
  \item
    \textbf{Validator scope.} A graph that passes identification is not
    proof of causality. Assumptions in \(\mathcal{A}\) stay untested.
  \end{itemize}

  \section{Conclusion and Future Work}\label{sec:conclusion}

  \begin{itemize}
  \tightlist
  \item
    Summary of contributions in 2 to 3 sentences.
  \item
    State plainly: design and protocol only, no implementation or
    empirical evaluation yet.
  \item
    Numbered next steps:

    \begin{enumerate}
    \def\labelenumi{\arabic{enumi}.}
    \tightlist
    \item
      Implement the Theory Spec schema.
    \item
      Build the tested library of estimators and data adapters.
    \item
      Run the anomaly-paper evaluation.
    \item
      Define the DAG optimization criterion.
    \item
      Build the literature-level graph.
    \end{enumerate}
  \end{itemize}

  \bibliographystyle{amsplain}
  \bibliography{references}
\end{document}
